%
% 5-argument.tex -- Argumentprinzip
%
% (c) 2026 Prof Dr Andreas Müller
%
\section{Das Argumentprinzip
\label{buch:integration:section:argumentprinzip}}
\kopfrechts{Das Argumentprinzip}%
Das Argument einer komplexen Zahl $z$ ist der Winkel $\varphi$,
mit dem $z=|z|\,e^{i\varphi}$ rekonstruiert werden kann.
Dieser ist nur bis auf ein Vielfaches von $2\pi$ bestimmt.
Diese Flexibilität wird benötigt, um die Entwicklung des
Arguments entlang einer Kurve zu untersuchen.
In Abschnitt~\ref{buch:integration:cauchy:subsetion:umlaufzahl}
wurde die Umlaufzahl als Integral mithilfe des Cauchy-Integralsatzes
berechnet.
In diesem Abschnitt soll jetzt gezeigt werden, dass das Bild
einer geschlossenen Kurve die Information darüber enthält, wieviele
Nullstellen und Polstellen eine meromorphe Funktion hat.

%
% 51-argument.tex
%
% (c) 2026 Prof Dr Andreas Müller
%
\subsection{Argumentänderungen}
Das Argument einer komplexen Zahl $z$ ist der Winkel $\varphi$, für den
$z=|z|e^{i\varphi}$.
Weil die imaginäre Exponentialfunktion wegen $e^{2\pi i}=1$ periodisch
ist, ist das Argument nur bis auf ein ganzzahliges Vielfaches von $2\pi$
festgelegt.
Dies wurde auch schon in
Abschnitt~\ref{buch:analytisch:fortsetzung:subsection:ueberlagerung}
beobachtet, wo die Logarthmusfunktion untersucht wurde.
Der Logarithmus der komplexen Zahl $re^{i\varphi}$ ist
\[
\log (re^{i\varphi}) = \log r + i\varphi + 2\pi i k
\]
für $k\in\mathbb{Z}$.
Die Logarithmusfunktion verspricht also, im Imaginärteil das
Argument zu liefern, versagt aber weil das Resultat nicht eindeutig
bestimmt ist.

%
% Das Argument einer Funktion
%
\subsubsection{Das Argument einer Funktion}
Für $z=0$ ist das Argument undefiniert, es hat bei $0$ eine
Singularität.
Für eine komplexe Funktion $f(z)$ ist es daher ausgeschlossen,
die Funktion $\log f(z)$ zu definieren, wenn die Funktion $f(z)$
eine Nullstelle hat.
Schliesst man Punkte des Definitionsbereichs aus, in denen $f(z)$
verschwindet, dann ist der Wert von $\log f(z)$ bis auf ein ganzzahliges
Vielfaches von $2\pi$ festgelegt.
Es ist also immer noch nicht möglich, einen eindeutigen Wert von $\log f(z)$
festzulegen.
Man kann aber einen der möglichen Werte für $\log f(z_0)$ an der Stelle
$z_0$ zu wählen und dann versuchen, die Werte von $\log f(z)$ für $z$ 
in einer Umgebung von $z_0$ so definieren, dass eine stetige
Funktion entsteht.

\begin{beispiel}
Die Funktion $f(z) = z^3$ bildet $z=re^{i\varphi}$ auf
$f(z) = r^3e^{3i\varphi}$ ab.
Der Logarithmus ist
\[
\log f(z)
=
\log r^3 + 3i\varphi
+
2\pi k
=
3\log r + 3i\varphi
+
2\pi i k
\]
für $k\in\mathbb{Z}$.
Um den Logarithmus festzulegen, muss man also den Wert für einen
Punkt  $z_0$ festzulegen.
Wir wählen $z_0=i=e^{i\frac{\pi}2}$, wir müssen also einen der Werte
\[
\dots
-\frac{9\pi}{2},
-\frac{5\pi}{2},
-\frac{\pi}{2},
{\color{darkred}\frac{3\pi}2}, 
\frac{7\pi}2,
\frac{11\pi}2,
\frac{15\pi}2,
\dots
\]
auswählen.
Der mittlere Wert $3\pi/2$ ist der offensichtliche, aber jede andere
Wahl ist genauso zulässig.

Mit der Wahl $\log f(i) = -\frac{\pi}2i$ können wir jetzt die anderen
Werte des Logarithmus in einer Umgebung von $z_0=i$ berechnen.
Auf dem Gebiet $\mathbb{C}\setminus\{x\in\mathbb{R}\mid x\le 0\}$,
der komplexen Ebene ohne die negative reelle Achse, kann das Argument
von $z=re^{i\varphi}$ im Intervall $\varphi\in(-\pi,\pi)$ eindeutig
festgelegt werden.
Mit dieser Wahl des Argumentes von $z$ ist dann
\[
\log f(z)
=
\log|f(z)| - \frac{\pi}2 + 3\biggl(\varphi-\frac{\pi}2\biggr)i.
\]
Mit dieser Wahl des Wertes von $\log f(z_0)$ haben positive reelle Werte
von $z$ das Argument $\varphi=0$ und damit die Werte
\[
\log f(r) = 3\log r - \frac{\pi}2 -3\frac{\pi}2 = 3\log r-2\pi
\]
des Logarithmus.
Wendet man die Exponentialfunktion an, ergibt sich der reelle Wert
$r^3$, da der Imaginärteil ein ganzzahliges Vielfaches von $2\pi$ ist.
\end{beispiel}

%
% Argument entlang einer Kurve
%
\subsubsection{Argument entlang einer Kurve}
Die Konstruktion von $\log f(z)$ im vorangegangenen Abschnitt
hat das Problem, dass der Logarithmus nicht eindeutig definiert werden kann,
in einer Umgebung eines einzelnen Wertes gelöst.
Verwendet wurde dabei, dass sich das Argument von $\log f(z)$ immer von
einem bekannten Wert $\log f(z_1)$ auf eine Umgebung ausdehnen lässt.
Ausgehend vom Wort in einem anderen Punkt $z_2$ dieser Umgebung, kann
$\log f(z)$ auf eine Umgebung von $z_2$ ausgedehnt werden.
Dies ist genau die Vorgehensweise der analytischen Fortsetzung einer
Funktion, die in Abschnitt~\ref{buch:analytisch:fortsetzung:subsection:kurve}
studiert wurde.
Sie garantiert, dass es entlang einer Kurve $\gamma$ im Definitionsbereich
der Funktion $f(z)$ immer möglich ist, den Logarithmus so zu wählen,
dass $\log f(\gamma(t))$ stetig ist.

Die Festlegung der Werte des Logarithmus von $f(z)$ entlang der Kurve
bedeutet, dass es stetige Funktionen $r(t)$ und $\varphi(t)$ mit
\begin{equation}
\log f(\gamma(t)) = \log r(t) + i\varphi(t)
\label{buch:analytisch:argument:argument:eqn:logf1}
\end{equation}
gibt, die
$f(\gamma(t))
=
e^{r(t)}\bigl(\cos\varphi(t)+i\sin\varphi(t)\bigr)$
erfüllen.

Für eine holomorphe Funktion $f(z)$ und eine differenzierbare Kurve
$\gamma(t)$ sind die Funktionen $r(t)$ und $\varphi(t)$ nicht nur
stetig, sondern sogar differenzierbar.
Man kann daher die Ableitung auf zwei verschiedene Arten berechnen, nämlich
einerseits mit der Kettenregel aus der Zusammensetzung von $\log\circ f$,
und andererseits aus \eqref{buch:analytisch:argument:argument:eqn:logf1}.
Die beiden Arten geben
\[
\frac{d}{dt}\log f(\gamma(t))
=
\begin{cases}
\displaystyle
\frac{f'(\gamma(t))}{f(\gamma(t))}
&\qquad\text{mit der Kettenregel}
\\[10pt]
\displaystyle
\frac{\dot{r}(t)}{r(t)} + i\dot{\varphi}(t)
&\qquad\text{aus 
\eqref{buch:analytisch:argument:argument:eqn:logf1}.}
\end{cases}
\]
Der Realteil beschreibt also die Änderung des Betrags $|f(z)|$ 
während der Imaginärteil die Änderung des Argumentes bestimmt.
Für eine auf dem Intervall $I=[a,b]$ definierte geschlossene Kurve
$\gamma$ kann man beide Formen der Ableitung entlang der Kurve
integrieren, esfolgt dann
\begin{align}
\oint_\gamma \frac{f'(z)}{f(z)}\,dz
&=
\int_I \frac{\dot{r}(t)}{r(t)}\,dt
+
i\int_I \dot{\varphi}(t)\,dt
\notag
\\
&=
\Bigl[ \log r(t) \Bigr]_a^b
+
i\Bigl[\varphi(t)\Bigr]_a^b
\notag
\\
&=
\bigl(\log r(b) - \log r(a)\bigr)
+
i
\bigl(\varphi(b)-\varphi(a)\bigr).
\label{buch:analytisch:argument:eqn:argumentintegral0}
\end{align}
Da die Kurve geschlossen ist, ist $r(b)=r(a)$, der Realteil von
\eqref{buch:analytisch:argument:eqn:argumentintegral0}
verschwindet daher.

\begin{satz}[Argumentintegral]
Für eine holomorphe Funktioin $f(z)$ und eine differenzierbare Kurve
$\gamma(t)$ im Definitionsgebiet von $f(z)$, auf der $f(z)$ nicht verschwindet,
ist das Integral
\begin{equation}
\oint_\gamma \frac{f'(z)}{f(z)}\,dz
=
2\pi i k
\label{buch:analytisch:argument:eqn:argumentintegral}
\end{equation}
mit $k\in\mathbb{Z}$.
\index{Argumentintegral}%
\end{satz}

Wir nennen \eqref{buch:analytisch:argument:eqn:argumentintegral}
das Argumentintegral.
Es hängt offenbar nicht von der konkreten Wahl des Logarithmus ab.
Insbesondere kann es auch ohne die Aufwendige konstruktion einer
analytischen Fortsetzung der Funktion $\log f(z)$ berechnet werden.



%
% 51-integral.tex -- Das Argumentintegral für f(z) = z^n
%
% (c) 2026 Prof Dr Andreas Müller
%
\subsection{Auswertung des Argumentintegrals
\label{buch:analytisch:argument:subsection:zn}}
In diesem Abschnitt soll die Verbindung zwischen dem Argumentintegral
\eqref{buch:analytisch:argument:eqn:argumentintegral0}
und dem früher eingeführten Begriff der Ordnung einer Nullstelle
oder eines Pols hergestellt werden.
Zu diesem Zweck wird zunächst die Funktion $f(z)=z^n$ 
betrachtet und später auf eine beliebige Nullstelle oder einen
Pol einer Funktion verallgemeinern.

%
% Geometrische Interpretation
%
\subsubsection{Geometrische Interpretation}
Das Argumentintegral
\eqref{buch:analytisch:argument:eqn:argumentintegral0}
berechnet die Änderung des Argumentes von $f(z)$, wenn man
einer geschlossenen Kurve folgt.
Für eine Funktion $f(z)$, die in $z=0$ weder einen Nullpunkt noch
einen Pol hat, und eine geschlossene Kurve $\gamma$, die den
Nullpunkte einmal in positiver Richtung einmal umläuft, ist
$f(\gamma(t))$ eine Kurve, entlang der sich das Argument um ein ganzzahliges
Vielfaches von $2\pi$ ändert.
Die ganze Zahl $k$, die in der Formel
\eqref{buch:analytisch:argument:eqn:argumentintegral} auftritt
ist also die Windungszahl des Weges $f(\gamma(t))$.

%
% Berechnung des Integrals
%
\subsubsection{Berechnung des Integrals für $f(z)=z^n$}
Wir betrachten zunächst die Funktionen $f(z)=z^n$ und berechnen 
das Integral
\begin{equation}
\frac{1}{2\pi i}
\oint_{\gamma}
\frac{f'(z)}{f(z)}
\,dz
=
\frac{1}{2\pi i}
\oint_{\gamma}
\frac{nz^{n-1}}{z^n}
\,dz
=
n\cdot
\frac{1}{2\pi i}
\oint_{\gamma}
\frac{1}{z}
\,dz
\label{buch:analytisch:argument:eqn:argumentintegral2}
\end{equation}
für eine Kurve $\gamma$, die den Nullpunkt nicht trifft und ihn
genau einmal im Gegenuhrzeigersinn umläuft.
Da $f'(z)/f(z)$ ausserhalb des Nullpunkts holomorph ist, können
wir den Weg auch durch den Einheitskreis $t\mapsto e^{2\pi it}$
ersetzen und werden den gleichen Wert für das Integral erhalten,
nämlich
\begin{align*}
\frac{1}{2\pi i}
\oint_{\gamma}
\frac{f'(z)}{f(z)}
\,dz
&=
n\cdot
\frac{1}{2\pi i}
\int_0^1
\frac{1}{e^{2\pi it}}
2\pi ie^{2\pi it}
\,dt
=
n
\int_0^1\,dt
=
n.
\end{align*}
Das Argumentintegral
\eqref{buch:analytisch:argument:eqn:argumentintegral2}
liefert also die Ordnung eines Poles oder einer Nullstelle.

%
% Argumentintegral für eine Funktion mit f(0)=0
%
\subsubsection{Argumentintegral für eine Funktion mit $f(0)=0$}
Die letzte Beobachtung soll noch etwas vertieft werden.
Wir betrachten eine holomorphe Funktion $f(z)$ mit einem
singulären Punkt endlicher Ordnung in $0$.
Für eine Nullstelle der Ordnung $n$ kann die Funktion $f$
als
\[
f(z)
=
z^n\, g(z)
\]
geschrieben, werden, wobei $g(z)$ eine holomorphe Funktion
ist, die im Nullpunkt nicht verschwindet.
Da Nullstellen isoliert sind, verschwindet $g(z)$ auch in
einer genügend kleinen Umgebung von $0$ nicht.

Die Ableitung
\[
f'(z)
=
nz^{n-1}g(z) + z^ng'(z).
\]
ist eine holomorphe Funktion.
Das Argumentintegral über einen kleine Kreis $\gamma$ um den
Nullpunkt ist
\begin{align*}
\frac{1}{2\pi i}
\oint_\gamma
\frac{f'(z)}{f(z)}
\,dz
&=
\frac{1}{2\pi i}
\oint_\gamma
\frac{nz^{n-1}g(z)+z^ng'(z)}{z^ng(z)}
\,dz
\\
&=
\frac{1}{2\pi i}
\oint_\gamma
\frac{n}{z}
+
\frac{g'(z)}{g(z)}
\,dz.
\intertext{Der zweite Term im Integral ist eine holomorphe Funktion
auch im Punkt $0$, der Weg $\gamma$ ist also in im Definitionsgebiet
von $g'(z)/g(z)$ zusammenziehbar und das Integral verschwindet.
Es bleibt daher}
&=
n\cdot
\frac{1}{2\pi i}
\oint_\gamma
\frac{1}{z}
\,dz
=
n,
\end{align*}
wie früher bereits ausgerechnet wurde.

%
% Argumentintegral für eine Polstelle
%
\subsubsection{Argumentintegral für eine Polstelle}
Hat die Funktion $f(z)$ einen Pol der Ordnung $n$, dann kann sie als
\[
f(z) = \frac{1}{z^n}g(z)
\]
geschrieben werden, wobei $g$ eine in einer Umgebung von $0$ holomorphe
Funktion ist, die keine Nullstellen oder Pole nahe bei $0$ hat.
Die Ableitung von $f$ ist
\[
f'(z)
=
-
\frac{n}{z^{n+1}}g(z)
+
\frac{1}{z^n}g'(z).
\]
Für diese Funktion ist das Argumentintegral
\begin{align*}
\frac{1}{2\pi i}
\oint_\gamma
\frac{f'(z)}{f(z)}
\,dz
&=
\frac{1}{2\pi i}
\oint_\gamma
\frac{
-
nz^{-(n+1)}g(z)
+
z^{-n}g'(z)
}{
z^{-n}\,g(z)
}
\,dz
\\
&=
\frac{1}{2\pi i}
\oint_\gamma
-
\frac{n}{z}
+
\frac{g'(z)}{g(z)}
\,dz.
\intertext{Wie im Fall einer Nullstelle ist der Weg $\gamma$
im Definitionsbereich des zweiten Terms Integranden zusammenziehbar,
der zweite Term trägt daher nichts zum Integral bei.
Für den ersten Term findet man}
&=
-n\cdot
\frac{1}{2\pi i}
\oint_\gamma
\frac{1}{z}
\,dz.
\end{align*}

%
% Argumentintegral und Ordnung
%
\subsubsection{Argumentintegral und Ordnung}
Die Berechnungen der vorangegangenen Abschnitte zeigen, dass
der Wert des Argumentintegrals für Nullstellen und Pole 
immer durch die Ordnung
\[
\omega(0;f)
=
\oint_{\gamma}
\frac{f'(z)}{f(z)}
\,dz
\]
ausgedrückt werden kann.
Wir haben damit den folgenden Satz bewiesen.

\begin{satz}[Ordnung]
\label{buch:analytisch:argument:satz:ordnung}
Sei $f$ eine holomorphe Funktion mit einer isolierten Singularität
an der Stelle $a\in\mathbb{C}$.
\index{Singularitat@Singularität}%
Sei $\gamma$ ein kleiner Kreis, der den Punkt $z_0$ einmal im
Gegenuhrzeigersinn umläuft.
Er sei zudem so klein, dass er keine anderen Null- oder Polstellen
von $f$ als $z_0$ enthält.
Dann gilt
\[
\omega(z_0;f)
=
\frac{1}{2\pi i}
\oint_{\gamma}
\frac{f'(z)}{f(z)}
\,dz.
\]
\end{satz}

%
% 52-meromorph.tex -- Das Argumentprinzip für beliebige meromorphe Funktionen
%
% (c) 2026 Prof Dr Andreas Müller
%
\subsection{Das Argumentprinzip für beliebige meromorphe Funktionen}
%
% fig-argument.tex -- 
%
% (c) 2026 Prof Dr Andreas Müller
%
\begin{figure}
\centering
\includegraphics{chapters/040-analytisch/images/argument.pdf}
\caption{Bild eines Kreises unter einer Funktion $f(z)$ in Abhängigkeit
von der Anzahl der Nullstellen (blaue Kreise) und Pole (blaue Kreuze),
die vom Kreis umlaufen werden.
Die Differenz ist die Anzahl Umläufe der Bildkurve um den Nullpunkt.
\label{buch:analytisch:fig:argument}}
\end{figure}
%
Der Satz~\ref{buch:analytisch:argument:satz:ordnung}
liefert Information für eine einzelne isolierte Singularität.
Wir möchten jetzt zeigen, dass sich die Schlussfolgerung auch
verallgemeinern lässt auf einen Weg, der mehr als eine Null-
oder Polstelle der Funktion $f$ umschliesst.

\begin{satz}[Argumentprintzip für meromorphe Funktionen]
\label{buch:analytisch:argument:satz:argument}
\index{Argumentprinzip!meromorphe Funktionen}%
\index{meromorph}%
Sei jetzt $f$ ein meromorphe Funktion und $\gamma$ ein einfach
geschlossener Weg in $\mathbb{C}$, der keine Nullstellen oder Polstellen
von $f$ trifft.
Dann gilt
\[
\frac{1}{2\pi i}
\oint_\gamma
\frac{f'(z)}{f(z)}
\,dz
=
\sum_{\text{$a$ Singularität von $f$}}
\omega(a;f),
\]
wobei die Summe über alle Null- oder Polstellen $a$ von $f$ 
zu erstrecken ist.
\end{satz}

\begin{proof}
Nach Satz~\ref{buch:integration:homologie:satz:pfadpunkte} ist das
Wegintegral über $\gamma$ gleich gross wie die Summe der Integrale
über kleine, im Gegenuhrzeigersinn durchlaufene Kreise um die einzelnen
Singularitäten.
Die Summanden haben nach
Satz~\ref{buch:analytisch:argument:satz:ordnung}
den Wert $\omega(a;f)$.
Damit ist der Satz bewiesen.
\end{proof}

In Abbildung~\ref{buch:analytisch:fig:argument}
wird das Bild eines Kreises unter einer Funktion $f(z)$ in Abhängigkeit
von der Anzahl der Nullstellen und Pole, die vom Kreis umlaufen werden.
Sie zeigt, dass die Bildkurve den Nullpunkt im Bildbereich soviele
Male umläuft, wie die Differenz der vom Kreis umlaufenen Nullstellen und
Pole angibt.

Der Satz \ref{buch:analytisch:argument:satz:argument} ermöglicht,
die Anzahl der Nullstellen und Polstellen einer Funktion zu zählen,
ohne dass man wissen muss, wo sich diese befinden.
Dies wird zum Beispiel im Algorithmus von Weyl zur numerischen
Bestimmung der komplexen Nullstellen eines Polynoms ausgenutzt, der in
Kapitel~\ref{chapter:weyl} beschrieben wird.



%
% 53-verallgemeinert.tex -- Das verallgemeinerte Argumentprinzip
%
% (c) 2026 Prof Dr Andreas Müller
%
\subsection{Das verallgemeinerte Argumentprinzip
\label{buch:analytisch:argument:subsection:verallgemeinert}}
Das Argumentprinzip von Satz~\ref{buch:analytisch:argument:satz:argument}
hat eine unmittelbare Verallgemeinerung.

\begin{satz}[Verallgemeinertes Argumentprinzip]
\index{Argumentprinzip!verallgemeinert}%
Sei $f$ eine meromorphe Funktion und $\gamma$ ein Weg, der keine Nullstellen
oder Pole von $f$ enthält und sei $g$ eine ganze Funktion.
Dann ist
\[
\oint_\gamma \frac{f'(z)}{f(z)} g(z)\,dz
=
\sum_{\text{$a$ Nullstelle von $f$}} \operatorname{ind}_\gamma(a) g(a)
-
\sum_{\text{$b$ Pol von $f$}} \operatorname{ind}_\gamma(b) g(a)
\]
Wobei die Nullstellen $a$ und Pole $b$ von $f$ mit Ihrer Vielfachheit
gezählt werden sollen.
\end{satz}


